% macroir_bessel_method.tex
% Draft theory/methods section for the MacroIR-Bessel paper.
% 2026-08-31. Source of record: theory/macroir/notes/macroir_bessel_plan.md
% Notation and vocabulary aligned with the submitted eLife manuscript
% (papers/1_method/docs/manuscript-drafts/): per-channel pair (mu, Sigma)
% with explicit N_ch factors as in 02_framework.tex Eq. block :543-551;
% "boxcar" / "acquisition kernel" as in Appendix 1 subsection "The
% acquisition kernel, and what the uniform window misses" (the canonical
% bound-table site); the augmented state (N, z) and its per-window reset as
% in Appendix 1 subsection on the augmented-state filter
% (08_appendix_derivation.tex:948-956). The object Sigma*gamma is named by
% its meaning (the covariance between the state a channel is in and the
% current it carries), never "gain".
% Status: exposition draft. The block-by-block derivation of the covariance
% prediction (pole-shifted endpoint-conditioned integrals, multinomial
% assembly) is milestone M0 of the plan; TODO markers flag what M0 fills.
% Claim discipline (plan, section 10): "extended closure, measured", never
% "exact likelihood"; the cost claim is stated as scaling, never as absence
% of prior art.

\documentclass[11pt]{article}

\usepackage{amsmath,amssymb,amsfonts,bm}
\usepackage{booktabs}
\usepackage{geometry}
\geometry{margin=1in}

\title{A likelihood that carries the recording filter's state\\[4pt]
Draft section for the MacroIR--Bessel paper}
\author{}
\date{}

\begin{document}

\maketitle

\newcommand{\E}{\mathbb{E}}
\newcommand{\Var}{\operatorname{Var}}
\newcommand{\Cov}{\operatorname{Cov}}
\newcommand{\diag}{\operatorname{diag}}
\newcommand{\Nch}{N_{\mathrm{ch}}}
\newcommand{\N}{\mathbf{N}}
\newcommand{\ff}{\bm{\varphi}}
\newcommand{\Af}{\mathbf{A}_{\mathrm{f}}}
\newcommand{\bff}{\mathbf{b}_{\mathrm{f}}}
\newcommand{\cf}{\mathbf{c}_{\mathrm{f}}}

%------------------------------------------------------------
\section{The measurement problem}
%------------------------------------------------------------

The patch-clamp amplifier low-pass filters the current before the digitizer
samples it, so the real acquisition kernel is the boxcar composed with the
amplifier's analogue low-pass. In the recordings that motivate this work
the filter is a 4-pole Bessel with cutoff $f_c = 10\,\mathrm{kHz}$ and the
sampling rate is $50\,\mathrm{kHz}$ \cite{moffatt2007}. A low-pass filter
has memory: the value the digitizer records at time $t$ is a weighted
average of the recent past of the current $i(t)$,
\begin{equation}
  z(t) \;=\; \int_0^{\infty} h(s)\, i(t-s)\,\mathrm{d}s
  \;+\; \text{filtered instrument noise},
  \label{eq:conv}
\end{equation}
where $h$ is the filter's impulse response. For a 4-pole Bessel at
$10\,\mathrm{kHz}$, $h$ is appreciable for about $100\,\mu\mathrm{s}$,
which spans five sampling intervals.

The likelihood of the companion paper \cite{macroir2026} implements the
boxcar exactly: each recorded value is modelled as the plain average of the
current over its own acquisition window. What the composition with the
analogue filter changes is governed by one dimensionless group, the cutoff
frequency against the window, $f_c\Delta$, and the companion paper's
appendix quantifies it: at $f_c\Delta = 0.2$ the white-noise component of a
recorded value has $61.5\%$ less variance than the boxcar assigns, and
neighbouring recorded values are correlated at $0.641$, both effects the
boxcar model ignores.

Two separate things go wrong when the kernel spills across window
boundaries:
\begin{enumerate}
  \item a recorded value depends on what the channels did during earlier
    windows, and the recursion never conditions on that history;
  \item consecutive recorded values share signal history and share
    filtered noise, so they are correlated even once the occupancies at
    the window boundaries are known, while the recursion treats them as
    conditionally independent.
\end{enumerate}

On the record of \cite{moffatt2007} the failure sits exactly where the
information is. The single-sample windows around the agonist pulses are
$20\,\mu\mathrm{s}$ long, so they sit at $f_c\Delta = 0.2$; the
$187\,\mu\mathrm{s}$ test pulses are only about $2.7$ times the filter's
fixed smear of $0.69/f_c = 69\,\mu\mathrm{s}$. Grouping samples until the
boxcar becomes valid would discard the very windows the ultrashort-pulse
experiment exists to produce. A filter-aware likelihood buys those windows
back.

The obstruction to fixing this inside the existing recursion is named in
the companion paper's appendix. There, the boundary-state construction is
shown equivalent to a filter on the augmented state $(\N, z)$, where $z$
accumulates the window average; the equivalence holds per window, with a
reset: the augmented filter must discard the posterior on $z$ and open the
next window with $z = 0$, because the next window's integral is fresh and
independent of the old given the occupancy. A filter that remembers across
windows forbids exactly that reset. The construction below removes the
need for it.

%------------------------------------------------------------
\section{The amplifier as part of the model}
%------------------------------------------------------------

A $k_{\mathrm{f}}$-pole filter is itself a small physical system. It has
$k_{\mathrm{f}}$ internal variables (the voltages on the capacitors inside
the filter), collected in a vector $\ff \in \mathbb{R}^{k_{\mathrm{f}}}$,
which evolve linearly, driven by whatever current enters:
\begin{equation}
  \frac{\mathrm{d}\ff}{\mathrm{d}t} \;=\; \Af\,\ff + \bff\, i(t),
  \qquad
  \text{recorded output} \;=\; \cf^{\top}\ff ,
  \label{eq:filter-ss}
\end{equation}
with $\Af,\bff,\cf$ fixed, known matrices determined by the pole positions
and the cutoff. For the 4-pole Bessel, $k_{\mathrm{f}} = 4$ and the
transfer function is $H(s) = \theta_4(0)/\theta_4(s/\omega_0)$ with
$\theta_4(s) = s^4 + 10s^3 + 45s^2 + 105s + 105$ the reverse Bessel
polynomial, realized in companion form, with the frequency scale
$\omega_0$ chosen so that the $-3\,\mathrm{dB}$ point falls at $f_c$. We
write $\nu_1,\dots,\nu_{k_{\mathrm{f}}}$ for the filter poles, with
$\operatorname{Re}\nu_p > 0$, so that the eigenvalues of $\Af$ are
$-\nu_p$; the poles come in complex-conjugate pairs and every quantity
below recombines to a real number.

The channels are the ensemble of the companion paper: $\Nch$ independent
Markov channels on $K$ states with generator $\mathbf{Q}$, transition
matrix $\mathbf{P}(t) = e^{\mathbf{Q}t}$ (row-stochastic), and conductance
vector $\bm{\gamma}$. Writing $\N(t) \in \mathbb{N}^K$ for the vector of
state counts, the current entering the amplifier is
\begin{equation}
  i(t) \;=\; \bm{\gamma}^{\top}\N(t) + i_{\mathrm{b}} + \xi(t),
  \label{eq:current}
\end{equation}
where $i_{\mathrm{b}}$ is the baseline current and $\xi$ is white
instrument noise with autocovariance
$\E[\xi(t)\,\xi(t')] = S_0\,\delta(t-t')$; $S_0$ is the noise power
spectral density referred to the filter input.

The digitizer does not record $i(t)$. It records the linear readout
$\cf^{\top}\ff(t)$ of the filter state at the sampling instants. This
observation converts the memory problem into a state problem. The recorded
value remembers more about the past than the occupancies at the window
boundaries do, which is why no recursion over occupancies alone can be
correct on short windows. All of that memory, however, is stored in $\ff$:
since $\ff$ is driven by $i$, the joint pair $(\N,\ff)$ is Markov, and its
present value determines the statistics of its whole future. Because the
filter sits in series with the channels, its states add to the description
($K + k_{\mathrm{f}}$ variables); this is unlike coloured noise modelled
as a parallel hidden process, which multiplies the state space.

%------------------------------------------------------------
\section{The augmented recursion}
%------------------------------------------------------------

\subsection{The carried state}

The companion paper's recursion carries a Gaussian in the per-channel pair
$(\bm{\mu},\bm{\Sigma})$, with $\bm{\mu} = \E[\N]/\Nch$ and
$\bm{\Sigma} = \Cov(\N)/\Nch$, conditioned on all recorded values so far.
The Bessel member carries a Gaussian over the joint of occupancies and
filter state, closed at two moments:
\begin{equation}
  \bm{\mu} = \frac{\E[\N]}{\Nch}, \quad
  \mathbf{m} = \E[\ff], \quad
  \bm{\Sigma} = \frac{\Cov(\N)}{\Nch}, \quad
  \bm{\Sigma}_{\!N\varphi} = \frac{\Cov(\N,\ff)}{\Nch}, \quad
  \bm{\Sigma}_{\!\varphi\varphi} = \Cov(\ff),
  \label{eq:carried}
\end{equation}
where the filter blocks $\mathbf{m}$ and
$\bm{\Sigma}_{\!\varphi\varphi}$ stay at physical scale because $\ff$ is
one physical object, driven by the whole ensemble.
$\bm{\Sigma}_{\!N\varphi}$ is the structurally new block: it records how
uncertainty about where the channels are is tied to uncertainty about the
filter's internal state, and it is the conduit through which a recorded
value informs the occupancy. Between windows nothing is reset: the
quantity whose reset the augmented formulation required, and a filter with
memory forbids, is simply carried.

\subsection{Prediction}

The posterior at the start of a window of length $\Delta$, pushed through
the joint dynamics, is the prior at its end. The means propagate through
the block-triangular augmented drift,
\begin{equation}
  \frac{\mathrm{d}}{\mathrm{d}t}
  \begin{pmatrix} \bm{\mu} \\ \mathbf{m} \end{pmatrix}
  =
  \begin{pmatrix}
    \mathbf{Q}^{\top} & \mathbf{0} \\
    \Nch\,\bff\,\bm{\gamma}^{\top} & \Af
  \end{pmatrix}
  \begin{pmatrix} \bm{\mu} \\ \mathbf{m} \end{pmatrix}
  +
  \begin{pmatrix} \mathbf{0} \\ \bff\, i_{\mathrm{b}} \end{pmatrix},
  \label{eq:mean-prop}
\end{equation}
so the spectrum of the augmented system is
$\operatorname{spec}(\mathbf{Q}) \cup \{-\nu_p\}$: the channel eigenvalues
and the filter poles, with no mixing. The covariance blocks propagate by
the laws of total expectation and variance over the hidden path,
conditioned on the boundary occupancies, exactly as in the companion
paper's boundary-state construction; the only new ingredients are the
endpoint-conditioned moments of exponentially weighted conductance
integrals, given in Section~\ref{sec:integrals}. Every expectation the
propagation needs is an affine or bilinear function of the moments already
carried in \eqref{eq:carried}, so the moment propagation itself introduces
no approximation.
% TODO(M0): the appendix derives each covariance block from the
% boundary-state construction (multinomial boundary counts, pole-shifted
% interval moments).

\subsection{The two reads}

The observation model splits into two cases, according to how the record
was produced.

\emph{Native-rate window} (one raw sample). The recorded value is the
filter output at the sampling instant,
\begin{equation}
  z \;=\; \cf^{\top}\ff(\Delta).
  \label{eq:read-native}
\end{equation}
There is no fresh observation noise in this read: the instrument noise
entered before the filter, and its contribution already lives inside
$\bm{\Sigma}_{\!\varphi\varphi}$ (Section~\ref{sec:noise}).

\emph{Grouped window} (an average of several raw samples). The recorded
value is the average of the filter output over the window. We add one
accumulator variable $u$, the direct heir of the augmented state $z$ of
the companion paper's appendix, now fed by the filter output rather than
by the raw current:
\begin{equation}
  \frac{\mathrm{d}u}{\mathrm{d}t} = \cf^{\top}\ff,
  \qquad u(0) = 0,
  \qquad
  z \;=\; \frac{u(\Delta)}{\Delta}.
  \label{eq:read-grouped}
\end{equation}
The reset of $u$ at the start of each recorded window is a property of the
recorder, whose averaging genuinely restarts there, so it remains as
legitimate as it was in the boxcar member. The accumulator is an
integrator, a pole at zero; it exists only within the window and is
discarded after the update. Applying the grouped read to a native-rate
sample would impose an averaging the hardware never performed, so the
member switches reads per window.

\subsection{Update}

The recorded value $z$ is a linear functional of the augmented state, so
the Bayes update is Gaussian conditioning of the joint moments, with the
same three blocks as the companion paper: the covariance of the state, the
variance of the recorded value, and the covariance between the state and
the recorded value. For the native-rate read,
\begin{align}
  z^{\mathrm{pred}} &= \cf^{\top}\mathbf{m},
  &
  \sigma^2 &= \cf^{\top}\bm{\Sigma}_{\!\varphi\varphi}\,\cf
             + \sigma_{\mathrm{fl}}^2,
  &
  \mathbf{g}_{\varphi} &= \bm{\Sigma}_{\!N\varphi}\,\cf,
  &
  \delta &= z - z^{\mathrm{pred}},
  \label{eq:threeblocks}
\end{align}
where $\mathbf{g}_{\varphi}$ takes the place of
$\mathbf{g} = \bm{\Sigma}\bm{\gamma}$ in the companion paper: it is the
covariance between the state a channel is in and the value the amplifier
is about to hand the digitizer. The term $\sigma_{\mathrm{fl}}^2$ is an
optional additive floor for noise components with no finite-dimensional
linear realisation (Section~\ref{sec:noise}). The posterior follows by
conditioning:
\begin{equation}
  \begin{aligned}
  \bm{\mu}^{\mathrm{post}}
    &= \bm{\mu} + \frac{\mathbf{g}_{\varphi}}{\sigma^2}\,\delta,
  &
  \bm{\Sigma}^{\mathrm{post}}
    &= \bm{\Sigma}
       - \frac{\Nch\,\mathbf{g}_{\varphi}\mathbf{g}_{\varphi}^{\top}}
              {\sigma^2},
  \\[4pt]
  \mathbf{m}^{\mathrm{post}}
    &= \mathbf{m}
       + \frac{\bm{\Sigma}_{\!\varphi\varphi}\cf}{\sigma^2}\,\delta,
  &
  \bm{\Sigma}_{\!\varphi\varphi}^{\mathrm{post}}
    &= \bm{\Sigma}_{\!\varphi\varphi}
       - \frac{(\bm{\Sigma}_{\!\varphi\varphi}\cf)
               (\bm{\Sigma}_{\!\varphi\varphi}\cf)^{\top}}{\sigma^2},
  \\[4pt]
  &
  &
  \bm{\Sigma}_{\!N\varphi}^{\mathrm{post}}
    &= \bm{\Sigma}_{\!N\varphi}
       - \frac{\mathbf{g}_{\varphi}
               (\bm{\Sigma}_{\!\varphi\varphi}\cf)^{\top}}{\sigma^2},
  \end{aligned}
  \label{eq:update}
\end{equation}
which reduces, block by block, to the companion paper's update when the
filter rows are removed. The grouped read conditions on $u(\Delta)/\Delta$
instead, with the same structure. Two things happen here that the boxcar
member cannot express. The datum updates the filter state as well as the
occupancy: it says where the capacitors are, and that knowledge carries
into the next window. And the correlation between neighbouring recorded
values is accounted for, because the predictive distribution of each
sample conditions on the previous ones through the carried filter moments.

\subsection{The likelihood}

The log-likelihood is the chain rule over windows,
\begin{equation}
  \log p(z_1,\dots,z_T)
  \;=\;
  \sum_{j=1}^{T} \log \mathcal{N}\!\bigl(z_j;\,
    z^{\mathrm{pred}}_j,\, \sigma^2_j\bigr),
  \label{eq:lik}
\end{equation}
with $z^{\mathrm{pred}}_j$ and $\sigma^2_j$ the predictive mean and
variance of sample $j$ given all earlier samples, read off
\eqref{eq:threeblocks} after the prediction step of window $j$.

%------------------------------------------------------------
\section{The interval integrals}
\label{sec:integrals}
%------------------------------------------------------------

\subsection{Pole decomposition}

The impulse response of the filter is a sum of decaying exponentials, one
per pole:
\begin{equation}
  h(s) \;=\; \cf^{\top} e^{\Af s}\, \bff
  \;=\; \sum_{p=1}^{k_{\mathrm{f}}} r_p\, e^{-\nu_p s},
  \label{eq:impulse}
\end{equation}
with residues $r_p$. Solving \eqref{eq:filter-ss} over one window,
\begin{equation}
  \ff(\Delta)
  \;=\;
  e^{\Af\Delta}\,\ff(0)
  \;+\;
  \int_0^{\Delta} e^{\Af(\Delta-s)}\,\bff\; i(s)\,\mathrm{d}s ,
  \label{eq:filter-sol}
\end{equation}
and in the eigenbasis of $\Af$ the channel part of the integral
decomposes, channel by channel, into the scalar quantities
\begin{equation}
  W_p \;=\; \int_0^{\Delta} e^{-\nu_p(\Delta-s)}\,
  \gamma_{X(s)}\,\mathrm{d}s ,
  \label{eq:Wp}
\end{equation}
where $X(s)$ is the state of one channel at time $s$. $W_p$ is the
channel's conductance history as the filter's $p$-th mode remembers it.
Setting $\nu_p = 0$ recovers the plain time integral of the conductance,
which is the boxcar member's object.

\subsection{Endpoint-conditioned moments and their spectral evaluation}

To propagate the covariance blocks, the recursion needs the moments of
$W_p$ conditioned on the channel's occupancy at the window boundaries, in
direct analogy with the boundary-conditioned window mean
$\bar{\gamma}_{i\to j}$ and per-state variance of the companion paper:
\begin{equation}
  \bar{\gamma}^{(p)}_{i\to j}
  = \E\!\left[ W_p \mid X(0)=i,\, X(\Delta)=j \right],
  \qquad
  \overline{M}^{(pq)}_{i}
  = \E\!\left[ W_p W_q \mid X(0)=i \right].
  \label{eq:cond-moments}
\end{equation}
The asymmetry in the conditioning is deliberate and it is where the cost
is contained. First moments are needed pair-resolved, because the
covariance between where a channel ends and what it contributed to the
filter is what feeds $\bm{\Sigma}_{\!N\varphi}$ and the update vector
$\mathbf{g}_{\varphi}$. Second moments enter only
$\bm{\Sigma}_{\!\varphi\varphi}$, where the end state is summed over, and
a pole-discounted propagator applied to the all-ones vector collapses,
\begin{equation}
  e^{(\mathbf{Q}-\nu\mathbf{I})t}\,\mathbf{1} = e^{-\nu t}\,\mathbf{1},
  \label{eq:collapse}
\end{equation}
so each pole pair $(p,q)$ costs a matrix-vector chain rather than a
$K\times K$ table.
% TODO(M0): prove the no-pair-resolved-seconds claim through every
% consumer, including the variance-corrected branches if inherited.

All of these integrals are evaluated by the machinery the boxcar member
already uses. With the spectral decomposition of the generator,
$\mathbf{Q} = \mathbf{V}\bm{\Lambda}\mathbf{V}^{-1}$ with eigenvalues
$\lambda_a$, the boxcar member's conditioned moments reduce to divided
differences of the exponential,
\begin{equation}
  D[x,y] = \frac{e^{x} - e^{y}}{x-y},
  \qquad D[x,x] = e^{x},
  \label{eq:divdiff}
\end{equation}
evaluated at pairs of arguments $\lambda_a\Delta,\ \lambda_b\Delta$. An
exponential weight in \eqref{eq:Wp} does nothing except shift arguments
inside the very same functions: the generic building block is
\begin{equation}
  \int_0^{\Delta} e^{\lambda_a s}\,
  e^{(\lambda_b-\nu_p)(\Delta-s)}\,\mathrm{d}s
  \;=\;
  \Delta\; D\!\left[\lambda_a\Delta,\ (\lambda_b-\nu_p)\Delta\right],
  \label{eq:shifted}
\end{equation}
and second moments involve the analogous second divided differences at
argument triples with shifts $\nu_p$, $\nu_q$. No new class of integral
appears; the existing code shapes run with shifted inputs, and setting
every $\nu_p = 0$ reproduces the boxcar member's formulas identically.

\subsection{Instrument noise}
\label{sec:noise}

The filter that colours the signal also colours the white instrument noise
$\xi$, and its contribution to the covariance of $\ff$ over a window has
the closed form
\begin{equation}
  \bm{\Sigma}_{\xi}(\Delta)
  \;=\;
  S_0 \int_0^{\Delta} e^{\Af s}\,\bff\,\bff^{\top} e^{\Af^{\top} s}\,
  \mathrm{d}s ,
  \label{eq:noise-gramian}
\end{equation}
which in the eigenbasis of $\Af$ is again the divided differences of
\eqref{eq:divdiff}, at arguments $-(\nu_p+\nu_q)\Delta$. This term
generalises, and replaces, the white per-sample instrumental variance
$\epsilon^2$ of the companion paper; the noise correlations it induces
across samples are then carried by $\bm{\Sigma}_{\!\varphi\varphi}$ like
any other filter memory. It also keeps the predictive variance $\sigma^2$
in \eqref{eq:threeblocks} strictly positive, since the pair $(\Af,\bff)$
is controllable. Noise components with no finite-dimensional linear
realisation, such as $1/f$ noise, stay outside the augmentation as the
additive per-observation floor $\sigma_{\mathrm{fl}}^2$.

%------------------------------------------------------------
\section{The boxcar member as an exact special case}
%------------------------------------------------------------

Delete the $\ff$ rows and keep only the accumulator $u$, so that the
observation system is $\Af = 0$, $\bff = 1$: the augmented state
collapses to the $(\N, z)$ of the companion paper's appendix, and the
equations of Section~\ref{sec:integrals} collapse, term by term, to the
boxcar member. The same collapse is visible in the algebra as the limit
$\nu_p = 0$ of \eqref{eq:shifted}. This subcase is the main regression
anchor of the implementation: the new code, run in that configuration,
must reproduce the existing, validated member to machine precision. The
limit $f_c \to \infty$ also recovers the boxcar member, though only
approximately at finite $f_c$, and serves as a secondary check.

%------------------------------------------------------------
\section{What is exact and what is approximate}
%------------------------------------------------------------

The moment propagation is exact: every expectation the recursion needs is
an affine or bilinear function of the moments it already carries, so no
information is silently dropped in the prediction step. What remains
approximate is exactly what is approximate in the boxcar member: the
Bayes update treats the joint distribution as Gaussian, and between
windows only two moments are carried. Both closures now act on a larger
object, the joint of occupancy and filter state, so their error is of the
same kind but need not be of the same size. We measure it on simulated
ground truth rather than assume it; the construction is an extended
closure whose accuracy is reported, and we make no exactness claim for
the likelihood itself.

%------------------------------------------------------------
\section{Cost}
%------------------------------------------------------------

Relative to the boxcar member, the growth is a combinatorial factor in
how many argument pairs the divided differences are evaluated at, with
the modes of the observation system playing the role of extra
eigenvalues: first-moment tables grow linearly in the number of modes
(one pair-resolved table per pole, plus one for the accumulator), while
second-moment quantities grow with the number of pole pairs but are held
at matrix-vector cost by \eqref{eq:collapse}. Conjugate symmetry halves
the complex arithmetic. The overall cost is polynomial in the state count
$K$ and the pole count $k_{\mathrm{f}}$, and independent of the number of
channels $\Nch$. Two properties of real records reduce it further:
distinct window lengths are few, so the window quantities are memoized on
the triple (length, stimulus, pole set); and on windows with
$f_c\Delta \gg 1$ the filter has settled well within the window, and the
computation reverts to the boxcar member's cost through an asymptotic
branch. On the record of \cite{moffatt2007}, 742 of 1078 windows sit in
that regime.

%------------------------------------------------------------
\section{Fixed filter parameters and identifiability}
%------------------------------------------------------------

The filter's cutoff and pole count are properties of the recording
apparatus. They are fixed at their known values and never fitted: they
trade almost degenerately against fast rate constants and against the
noise amplitude, so releasing them would launder model error into
pseudo-physical parameters. Rates faster than $f_c$ remain weakly
identified under any correct treatment, because the filter genuinely
destroyed that information before the digitizer saw it; the honest
outcome is wider posteriors on those rates, and a narrower posterior
there would be a defect of the method, a point the simulations verify.
Under a filter the noise amplitude must be stated as the input-referred
spectral density $S_0$ of \eqref{eq:current}; conventions that quote a
variance per sample depend on the bandwidth and do not transfer across
members.
% TODO: priors on S_0 restated accordingly before any evidence comparison
% across members (plan, decision 4).

%------------------------------------------------------------
\section{Relation to existing treatments}
%------------------------------------------------------------

Filters inside single-channel likelihoods have a long history. Michalek
and colleagues identified the anti-aliasing filter as a cause of wrong
model identification and modelled it with moving-average filtered hidden
Markov models \cite{michalek1999}; Venkataramanan, Kuc and Sigworth built
the metastate formulation, expanding the hidden state over the filter's
memory so that band-limited sampled data become Markov again
\cite{venkataramanan2000}; Qin, Auerbach and Sachs modelled the filter as
a finite impulse response within a hidden Markov likelihood and describe
an extension to multiple channels \cite{qin2000}. These treatments are
discrete-time and truncate the filter to a sampled impulse response, with
a cost that grows exponentially in the filter memory expressed in states
per channel, which confines them to one or few channels. Absorbing a
linear filter into an augmented state is textbook filtering theory
\cite{jazwinski1970}; propagating integrated or averaged observations
through an augmented state with matrix-exponential machinery goes back at
least to Zadrozny \cite{zadrozny1988}, has a general integral primitive
in Carbonell, Jim\'enez and Pedroso \cite{carbonell2008}, and a recent
instance with overlapping windows in Rubenzahl and colleagues
\cite{rubenzahl2026}.

What this member contributes is the combination needed for macroscopic
channel data: the endpoint-conditioned mean, variance and cross-sample
covariance of a Bessel-filtered current for $\Nch$ independent channels,
computed from the generator's spectral decomposition at a cost polynomial
in state count and pole count and independent of $\Nch$, exact for the
analogue kernel where the treatments above are exact only for a sampled
truncation; and, with it, the ability to fit a macroscopic record at its
native sampling rate, including the windows the filter smears.

%------------------------------------------------------------
\begin{thebibliography}{9}

\bibitem{moffatt2007}
Moffatt, L. \& Hume, R.~I. (2007).
Responses of rat P2X2 receptors to ultrashort pulses of ATP provide
insights into ATP binding and channel gating.
\emph{Journal of General Physiology} 130(2):183--201.

\bibitem{macroir2026}
% TODO: replace with the final citation of the companion (MacroIR) paper.
Moffatt, L. (2026). The companion MacroIR paper (submitted).

\bibitem{michalek1999}
Michalek, S., Lerche, H., Wagner, M., Mitrovi\'c, N., Schiebe, M.,
Lehmann-Horn, F. \& Timmer, J. (1999).
On identification of Na$^+$ channel gating schemes using moving-average
filtered hidden Markov models.
\emph{European Biophysics Journal} 28(7):605--609.

\bibitem{venkataramanan2000}
Venkataramanan, L., Kuc, R. \& Sigworth, F.~J. (2000).
Identification of hidden Markov models for ion channel currents.
III. Bandlimited, sampled data.
\emph{IEEE Transactions on Signal Processing} 48(2):376--385.

\bibitem{qin2000}
Qin, F., Auerbach, A. \& Sachs, F. (2000).
Hidden Markov modeling for single channel kinetics with filtering and
correlated noise.
\emph{Biophysical Journal} 79(4):1928--1944.

\bibitem{jazwinski1970}
Jazwinski, A.~H. (1970).
\emph{Stochastic Processes and Filtering Theory.}
Academic Press.

\bibitem{zadrozny1988}
Zadrozny, P. (1988).
Gaussian likelihood of continuous-time ARMAX models when data are stocks
and flows at different frequencies.
\emph{Econometric Theory} 4(1):108--124.

\bibitem{carbonell2008}
Carbonell, F., Jim\'enez, J.~C. \& Pedroso, L.~M. (2008).
Computing multiple integrals involving matrix exponentials.
\emph{Journal of Computational and Applied Mathematics} 213(1):300--305.

\bibitem{rubenzahl2026}
Rubenzahl, R.~A., Hattori, S., S\"arkk\"a, S., Farr, B.~F., Luhn, J.,
Foreman-Mackey, D. \& Bedell, M. (2026).
Scalable Gaussian processes for integrated and overlapping measurements
via augmented state-space models.
\emph{Astronomical Journal}, doi:10.3847/1538-3881/ae4d0b.

\end{thebibliography}

\end{document}
